\documentclass[11pt]{article}
\usepackage[margin=1in]{geometry}
\usepackage{amsmath,amssymb,amsthm,hyperref}
\newtheorem{theorem}{Theorem}
\newtheorem{lemma}{Lemma}
\newcommand{\E}{\mathbb E}
\title{All alternating ratio inequalities for the\\
Fennessey--Larcombe--French sequence}
\author{Principia Math In Progress}
\date{11 September 2026}
\begin{document}
\maketitle
\begin{abstract}
We prove Conjecture 5.3 of Sun and Zhao for every ratio iteration and
every required index, with strict inequalities. The normalized sequence
is a constant minus moments of an explicit positive measure on $(0,1)$.
A convergent logarithmic series then expresses its logarithmic deficit
as positive moments, making all higher finite-difference signs explicit.
\end{abstract}

\section*{Task 1: Resolve MathDB problem 332484}
Let $V_0=1$, $V_1=8$, and, for $n\geq2$, define
\begin{equation}\label{eq:Vrec}
 (n-1)n^2V_n
 =8(n-1)(3n^2-n-1)V_{n-1}-128(n-2)n^2V_{n-2}.
\end{equation}
For a positive sequence $a$, set $(\mathcal Ra)_n=a_{n+1}/a_n$ and
$(\Delta f)_n=f_{n+1}-f_n$.
The conjecture \cite[Conjecture 5.3, p.~601]{published} asserts that
$\mathcal R^k(V_n)_{n\geq1}$ is log-convex for odd $k\geq1$ and
log-concave for even $k\geq1$.

\begin{theorem}\label{thm:main}
For every $n\geq1$ and integer $r\geq2$,
\begin{equation}\label{eq:sign}
 (-1)^{r+1}\Delta^r\log V_n>0.
\end{equation}
Consequently, every inequality in Conjecture 5.3 holds strictly.
\end{theorem}

\subsection*{Identifying a moment representation}
Use the parameter convention for the complete elliptic integral:
\[
 K(m)=\int_0^{\pi/2}\frac{d\theta}{\sqrt{1-m\sin^2\theta}},\qquad
 c_j=\left(\frac{\binom{2j}j}{4^j}\right)^2.
\]
The binomial expansion under the integral gives
\begin{equation}\label{eq:Kseries}
 K(m)=\frac\pi2\sum_{j\geq0}c_jm^j\qquad(0\leq m<1).
\end{equation}
Put $L=2/\pi$ and, for $0<u<1$, $u\ne1/2$, define
\begin{align}
 w(u)&=\frac4{\pi^2}K\bigl(4u(1-u)\bigr),\label{eq:w}\\
 \rho(u)&=w(u)+(2u-1)w'(u).
 \label{eq:rho}
\end{align}
The value at $u=1/2$ is immaterial for integration.

\begin{lemma}\label{lem:moment}
The function $\rho$ is positive almost everywhere on $(0,1)$, integrable,
and satisfies
\begin{equation}\label{eq:mass}
 \int_0^1\rho(u)\,du=\frac4\pi-1=2L-1.
\end{equation}
For every $n\geq1$,
\begin{equation}\label{eq:moment}
 \frac{V_n}{16^n}=L-\int_0^1u^n\rho(u)\,du.
\end{equation}
\end{lemma}

We first identify the sequence that occurs in the integral. Let $X,Y$
be independent random variables with arcsine density
$1/(\pi\sqrt{1-x^2})$ on $(-1,1)$, and put $U=(1+XY)/2$.
A change of variables in the product density gives, for $0<|t|<1$,
\[
 f_{XY}(t)=\frac2{\pi^2}\int_{|t|}^1
 \frac{dx}{\sqrt{(1-x^2)(x^2-t^2)}}
 =\frac2{\pi^2}K(1-t^2).
\]
For the second equality set
$x^2=t^2+(1-t^2)\sin^2\theta$ and reverse the angle.
Thus $U$ has probability density $w$ in \eqref{eq:w}.
Since $\E(X^{2j})=\binom{2j}j/4^j$ and odd moments vanish,
\begin{equation}\label{eq:pn}
 p_n:=\E(U^n)=2^{-n}\sum_{j=0}^{\lfloor n/2\rfloor}\binom n{2j}c_j.
\end{equation}

For completeness, we verify from these finite sums that the $V_n$
in \eqref{eq:Vrec} are exactly the sequence needed below.
Let $H(q)=\sum_{j\geq0}c_jq^j$. The ratio of consecutive $c_j$ gives
\[
 q(1-q)H''(q)+(1-2q)H'(q)-\tfrac14H(q)=0.
\]
Summing \eqref{eq:pn} for $|z|$ small yields
\[
 A(z):=\sum_{n\geq0}p_nz^n
   =\frac2{2-z}H\left(\frac{z^2}{(2-z)^2}\right).
\]
Substitution into the preceding differential equation, by the chain
rule, gives
\[
 z(1-z)(2-z)A''+(2-6z+3z^2)A'+(z-1)A=0.
\]
Its coefficient equations show that $P_n:=16^np_n$ satisfies
\[
 E_n:=n^2P_n-8(3n^2-3n+1)P_{n-1}+128(n-1)^2P_{n-2}=0
 \quad(n\geq2),
\]
with $P_0=1$, $P_1=8$; also set $E_1=P_1-8P_0=0$.
Define $W_0=1$ and $W_n=(2n+1)P_n-16nP_{n-1}$ for $n\geq1$.
The residual of \eqref{eq:Vrec} with $W$ in place of $V$ is exactly
\[
 (n-1)(2n+1)E_n-16n^2E_{n-1}\qquad(n\geq2).
\]
It is therefore zero. Since $W_1=8$, uniqueness in the recurrence gives
$W_n=V_n$, and hence
\begin{equation}\label{eq:weighted}
 \frac{V_n}{16^n}=(2n+1)p_n-np_{n-1}
 =\int_0^1\left[u^n+n(2u-1)u^{n-1}\right]w(u)\,du
 \quad(n\geq1).
\end{equation}

\begin{proof}[Proof of Lemma \ref{lem:moment}]
Write $m=4u(1-u)=1-(2u-1)^2$. Differentiating \eqref{eq:w} gives
\[
 \rho(u)=\frac4{\pi^2}\left[K(m)-4(1-m)K'(m)\right].
\]
Here $K'$ differentiates its parameter $m$.
Using \eqref{eq:Kseries} and
$c_{j+1}/c_j=(2j+1)^2/(4(j+1)^2)$, we obtain the positive series
\begin{equation}\label{eq:positive}
 \rho(u)=\frac2\pi\sum_{j\geq1}\frac{j}{j+1}c_j
                         \bigl(4u(1-u)\bigr)^j>0
 \qquad(0<u<1,\ u\ne1/2).
\end{equation}

We justify the integration by parts even at the middle singularity.
The elementary bound $c_j\leq1/(j+1)$ follows inductively from the
displayed coefficient ratio. Thus \eqref{eq:Kseries} implies
$K(m)=O(1+|\log(1-m)|)$ as $m\uparrow1$.
It follows that $h(u):=(2u-1)w(u)$ extends continuously across $1/2$
with value zero. At the endpoints, $w(0)=w(1)=L$, so $h(0)=-L$ and
$h(1)=L$. Away from the middle point,
\[
 h'(u)=w(u)+\rho(u)>0.
\]
Integrating separately on either side and taking the one-sided limits
shows that $h'$ is integrable with integral $2L$. Since $w$ is a
probability density, this proves \eqref{eq:mass} and the integrability
of $\rho$. It also proves that the continuous extension of $h$ is
absolutely continuous on $[0,1]$.

For $n\geq1$, integrate the second term of \eqref{eq:weighted} by parts:
\[
 \int_0^1\left[u^nw(u)+h(u)(u^n)'\right]du
 =\left[u^nh(u)\right]_0^1
       +\int_0^1u^n\bigl(w(u)-h'(u)\bigr)du
 =L-\int_0^1u^n\rho(u)du.
\]
This is \eqref{eq:moment}. The boundary term at zero vanishes precisely
because $n\geq1$; the formula does not assert a representation of $V_0$.
\end{proof}

\subsection*{The logarithm as positive moments}
Let $\mu$ be the positive measure with density $\rho(u)/L$ on $(0,1)$.
Its total mass is
\[
 q:=\mu((0,1))=2-\frac\pi2\in(0,1).
\]
For each positive integer $j$, let $\mu^{\odot j}$ be the image of
$\mu^{\otimes j}$ under multiplication
$(u_1,\ldots,u_j)\mapsto u_1\cdots u_j$. The measure
\begin{equation}\label{eq:tau}
 \tau=\sum_{j\geq1}\frac1j\mu^{\odot j}
\end{equation}
is finite and positive, since its mass is $\sum_{j\geq1}q^j/j<\infty$.
For $n\geq1$, Lemma \ref{lem:moment}, the power series for $-\log(1-z)$,
and Tonelli's theorem now give
\begin{align}
 \log V_n
 &=n\log16+\log L-\int_{(0,1)}t^n\,d\tau(t).
 \label{eq:logmoment}
\end{align}
In particular all $V_n$ used here are positive, since
$V_n/16^n\geq L(1-q)>0$.

\begin{proof}[Proof of Theorem \ref{thm:main}]
For $r\geq2$, the affine terms in \eqref{eq:logmoment} disappear.
Taking a finite difference inside the integral is legitimate because
$\tau$ is finite. Therefore
\[
 (-1)^{r+1}\Delta^r\log V_n
   =\int_{(0,1)}t^n(1-t)^r\,d\tau(t)>0.
\]
Strictness follows already from the first summand $\mu$ in
\eqref{eq:tau}, whose density is positive almost everywhere in $(0,1)$.
Finally, $\log(\mathcal R^kV)_n=\Delta^k\log V_n$, so
\[
 \Delta^2\log(\mathcal R^kV)_n=\Delta^{k+2}\log V_n.
\]
This is positive for odd $k$ and negative for even $k$, exactly the
required strict log-convexity and strict log-concavity for all $n\geq1$.
\end{proof}

\subsection*{Scope and verification}
The theorem covers every iteration order with the exact initial index
in the published conjecture. No asymptotic cutoff in $n$ is used.
The exceptional value $V_0$ is deliberately excluded, as in the source.
The proof supplies the full argument; finite symbolic and numerical
checks are supplementary. The accompanying audit records the current
published source, live MathDB status, and bounded search for solution
claims. No external peer review is claimed.

\section*{Task 2: Resolve the determinant-sequence conjecture}
Use the positive density $\rho$ to analyze
$D_n=V_n^2-V_{n-1}V_{n+1}$ for $n\geq2$.
Resolve published Conjecture 5.1, which asserts infinite log-convexity
of this sequence (MathDB 332482). In particular, determine whether
an appropriate exponential normalization of $D_n$ admits a positive
moment measure and identify that measure if it exists.

\section*{Task 3: Resolve the companion Catalan ratio conjecture}
For the Catalan--Larcombe--French numbers $P_n$ identified above,
resolve published Conjecture 5.2 (MathDB 332483): after omitting the
first $k$ terms, $\mathcal R^k(P_n)_{n\geq0}$ is log-concave for odd
$k$ and log-convex for even $k$. Explain how the necessary index cutoff
arises from the arcsine-product moment representation.

\section*{Task 4: Extend the construction to symmetric beta laws}
Replace the two arcsine densities by densities proportional to
$(1-x^2)^{a-1}$ and $(1-y^2)^{b-1}$ on $(-1,1)$, with $a,b>0$.
For $U=(1+XY)/2$ and
$B_n=(2n+1)\E(U^n)-n\E(U^{n-1})$, determine the parameter region
where $B_n$ is positive and has the same higher logarithmic-difference
signs. Account for any endpoint terms that replace the constant $L$.

\section*{Task 5: Obtain uniform quantitative inequalities}
Find explicit upper and lower bounds for
$(-1)^{r+1}\Delta^r\log V_n$ using the measure $\tau$.
Derive sharp asymptotics both for fixed $r$ with $n\to\infty$ and when
$r$ grows proportionally to $n$, and identify the contributions of
the individual multiplicative-convolution terms.

\clearpage
\begin{thebibliography}{9}
\bibitem{published}
B. Y. Sun and J. J. Y. Zhao, \emph{Log-behavior of Two Sequences Related
to the Elliptic Integrals}, Acta Mathematicae Applicatae Sinica,
English Series \textbf{36} (2020), 590--602.
Conjecture 5.3, p.~601.
\url{https://doi.org/10.1007/s10255-020-0952-1}.
Full publisher PDF:
\url{https://applmath.cjoe.ac.cn/jweb_yysxxb_en/EN/PDF/10.1007/s10255-020-0952-1}.
\bibitem{preprint}
B. Y. Sun and J. J. Y. Zhao, \emph{Log-behavior of two sequences related
to the elliptic integrals}, arXiv:1602.04359v4 (2017), Conjecture 5.3,
p.~14. \url{https://arxiv.org/abs/1602.04359}.
\end{thebibliography}
\end{document}
